\chapter*{摘\quad 要}
\addcontentsline{toc}{chapter}{摘要}

随着人工智能、人机交互和移动互联网技术的发展，面向听障群体的辅助交流与训练系统逐渐成为无障碍信息技术的重要应用方向。本文围绕听障用户在日常交流、手语学习、口语训练和管理端资源维护等场景中的实际需求，设计并实现了一套名为“听桥”的多端协同系统。系统由 HarmonyOS 移动端、Spring Boot 后端、Python 手势识别服务和 Web 管理端组成，支持手语资源展示、数字人手语播放、实时手势识别、训练样本采集、模型训练管理和模型版本发布等功能。

本文首先分析了系统的研究背景、用户需求和功能目标；随后设计了移动端、服务端、识别服务和管理端协同工作的总体架构；在实现层面，移动端负责用户交互、相机采集和识别结果展示，后端负责业务数据管理与接口服务，Python 服务负责基于 MediaPipe 和深度学习模型的手势识别处理，Web 管理端负责资源、样本、训练任务和模型版本的管理。最后，本文对系统功能和关键流程进行了测试，验证了系统在手势识别、样本采集、模型发布和多端协同方面的可行性。

\vspace{1em}

\noindent\textbf{关键词：}听桥；听障辅助；手语识别；HarmonyOS；Spring Boot；深度学习
